\chapter{Nebensätze}
\label{cha:Nebensaetze}

Nebensätze ergänzen oder erweitern den Hauptsatz. Sie werden durch Konjunktionen eingeleitet.

\section{Konjunktionen mit Verb am Ende}
\label{sec:Konjunktionen}

\subsection{Kausal (Ursache)}
\begin{itemize}
    \item \textbf{weil}: ..., weil sie krank war. (Nebensatz)
    \item \textbf{da}: Da ich müde war, ging ich ins Bett.
    \item \textbf{denn}: ..., denn ich war müde. (Hauptsatz!)
\end{itemize}

\subsection{Konzessiv (Gegensatz)}
\begin{itemize}
    \item \textbf{obwohl}: ..., obwohl es kalt war.
    \item \textbf{trotzdem}: ..., sie arbeitete trotzdem weiter. (Adverb, Hauptsatz!)
\end{itemize}

\subsection{Temporal (Zeit)}
\begin{itemize}
    \item \textbf{bevor}: ..., bevor ich ging.
    \item \textbf{nachdem}: ..., nachdem ich gegessen hatte.
    \item \textbf{als}: ..., als ich Kind war.
    \item \textbf{während}: ..., während ich las.
    \item \textbf{sobald}: ..., sobald ich kann.
    \item \textbf{bis}: ..., bis ich fertig bin.
\end{itemize}

\subsection{Konditional (Bedingung)}
\begin{itemize}
    \item \textbf{wenn/falls}: ..., wenn ich Zeit habe.
    \item \textbf{ohne dass}: ..., ohne dass ich wusste.
    \item \textbf{um zu}: ..., um zu lernen.
\end{itemize}

\subsection{Final (Zweck)}
\begin{itemize}
    \item \textbf{damit}: ..., damit du verstehst.
    \item \textbf{um zu}: ..., um zu lernen.
\end{itemize}

\section{Übungen zu Konjunktionen}
\label{sec:KonjunktionenUebungen}

\subsection{Kausal}
\begin{enumerate}
    \item Sie hat den Bus verpasst. Sie kam zu spät. $\rightarrow$ weil \quad \textit{(Lösung: Sie kam zu spät, weil sie den Bus verpasst hatte.)}
    \item Ich bleibe zu Hause. Ich bin krank. $\rightarrow$ denn \quad \textit{(Lösung: Ich bin krank, denn ich bleibe zu Hause.)}
    \item Es ist kalt. Wir bleiben drinnen. $\rightarrow$ da \quad \textit{(Lösung: Da es kalt ist, bleiben wir drinnen.)}
    \item Er hat nicht gelernt. Er hat die Prüfung nicht bestanden. $\rightarrow$ weil \quad \textit{(Lösung: Er hat die Prüfung nicht bestanden, weil er nicht gelernt hatte.)}
\end{enumerate}

\subsection{Konzessiv}
\begin{enumerate}
    \item Es war kalt. Wir sind schwimmen gegangen. $\rightarrow$ obwohl \quad \textit{(Lösung: Obwohl es kalt war, sind wir schwimmen gegangen.)}
    \item Er ist arm, aber er ist glücklich. $\rightarrow$ obwohl \quad \textit{(Lösung: Obwohl er arm ist, ist er glücklich.)}
    \item Sie hatte Kopfschmerzen. Sie arbeitete weiter. $\rightarrow$ trotzdem \quad \textit{(Lösung: Sie hatte Kopfschmerzen; sie arbeitete trotzdem weiter.)}
    \item Der Film war langweilig. Wir blieben bis zum Ende. $\rightarrow$ trotzdem \quad \textit{(Lösung: Der Film war langweilig; wir blieben trotzdem bis zum Ende.)}
\end{enumerate}

\chapterend